\section{Trabalhos Relacionados}
\label{sec:related}

A literatura prévia corre por quatro fios: a medição do viés factual
geográfico, extensões multilíngues e culturais, tentativas de mitigação, e o uso
de LLMs em contextos de decisão de política e saúde.

\subsection{Viés factual geográfico e extensões multilíngues}

\citet{manvi2024llm} estabeleceram que LLMs de fronteira carregam vieses
sistemáticos contra regiões de condições socioeconômicas mais baixas, com
correlações de até 0,70 entre avaliações geradas pelo modelo e um proxy dessas
condições; o espaço de desfecho são avaliações subjetivas, o que deixa em aberto
se o mesmo viés aparece na recordação objetiva de fatos regulatórios.
\citet{moayeri2024worldbench} quantificaram a recordação factual contra
indicadores do Banco Mundial em 20 LLMs, documentando erro relativo absoluto
$1{,}5\times$ maior para a África Subsaariana do que para a América do Norte; o
desenho é limitado pelo espaço de indicadores macroeconômicos, enquanto um
gestor ambiental faz perguntas mais estreitas e específicas de fonte, cujo
gabarito vive em conselhos ambientais e agências de monitoramento.
\citet{mirza2024global} mostraram que versões mais novas do GPT-4 não melhoram
monotonicamente e que o privilégio de afirmações sobre o Norte Global persiste.

A literatura de sondagem multilíngue estabeleceu cedo que o que um modelo
recorda depende do idioma da consulta. \citet{jiang2020xfactr} sondaram fatos
em formato cloze em 23 idiomas e encontraram acurácia abaixo de 15\% mesmo em
inglês e espanhol e abaixo de 5\% em marathi e iorubá, com metade dos fatos
corretamente recordados recordados em um único idioma; \citet{kassner2021mlama}
traduziram os mesmos benchmarks para 53 idiomas e encontraram o mBERT abaixo
de 60\% do seu desempenho em inglês em 32 deles, um viés para entidades do
idioma da consulta e nenhuma relação clara entre o tamanho da Wikipédia de um
idioma e sua recordação. Esses resultados dizem respeito a codificadores
mascarados de 2020; a pergunta que levantam, se a recordação em um idioma
acompanha o corpus desse idioma, é a que H2 e H4 fazem a modelos generativos.
\citet{myung2024blend} cobriram 16 países e 13 idiomas em conhecimento
cultural cotidiano, constatando que os LLMs têm melhor desempenho
nos idiomas nativos de culturas de alto recurso e pior nos de baixo recurso, o
achado que motiva H2 e se conecta à hierarquia de recursos de
\citet{joshi2020state}. Benchmarks regionais proliferaram:
BRoverbs~\citep{broverbs2025}, TiEBe~\citep{almeida2025tiebe},
ALBA~\citep{alba2026} e AMALIA~\citep{amalia2026} documentam falhas concretas
em tarefas culturalmente situadas, mas não tratam do uso aplicado em política
pública. O trabalho de construção de corpora em português é explícito em que o
esforço se concentrou no inglês: \citet{almeida2025portuguese} montam um corpus
de 120 bilhões de tokens em português porque pipelines específicos por idioma
são necessários para alcançar qualidade industrial. Essa assimetria na
quantidade de texto por idioma é o que H4 testa.

Esses estudos estabelecem que o fenômeno existe e é mensurável. Nenhum isola um
único domínio aplicado com gabarito oficial verificável, e nenhum trata o
enquadramento do prompt como fator manipulável.

\subsection{Mitigação, persona e condicionamento por papel}

\citet{opuszko2026unraveling} constataram que o raciocínio ajuda no agregado,
mas não elimina as disparidades regionais, e \citet{kerche2026silicon}
propuseram uma tipologia baseada em lugar que situa o viés geográfico num
panorama sociotécnico mais amplo. Ambos são conceituais e agregados.

O condicionamento por papel é o que os fornecedores instruem os próprios
usuários a fazer. A Anthropic documenta a definição de papel no prompt de
sistema~\citep{anthropic2026systemprompts}, o guia da OpenAI abre a estrutura
de mensagem recomendada com um bloco de identidade~\citep{openai2026promptengineering},
a documentação do Gemini, da Google, pede que definições de papel, que ela
chama de persona, sejam colocadas na instrução de
sistema~\citep{google2026promptingstrategies}, o exemplo padrão da xAI inicia
toda conversa com um prompt de sistema que atribui um
papel~\citep{xai2026chatguide}, e o padrão está catalogado na literatura de
prompting~\citep{white2023prompt}. A evidência sistemática sobre se isso ajuda
é negativa: em quatro famílias de LLM e 2.410 perguntas factuais, personas de
especialista não melhoraram a acurácia em relação a um controle sem persona e
a reduziram ligeiramente em 162 papéis~\citep{zheng2024helpful}. Um estudo
independente com seis modelos em questões de múltipla escolha de nível de
pós-graduação chega à mesma conclusão: casar uma persona de especialista com o
tipo de problema não teve efeito significativo no
desempenho~\citep{mollick2025personas}. Se o condicionamento por persona ajuda
ou prejudica a confiabilidade factual em domínios de alto risco é, portanto,
questão em aberto e, até onde sabemos, não testada para equidade geográfica.
Tratamos a persona como fator experimental pré-especificado e testamos tanto a
hipótese de conformidade à norma (a persona reduz a lacuna) quanto a hipótese
do vácuo de persona (a persona a amplia).

\subsection{LLMs em política, saúde e na administração}

\citet{lecoz2025policymaking} avaliaram recomendações de política de LLMs
para a situação de rua em quatro cidades (Barcelona, Joanesburgo, South Bend e
Macau) e leram o padrão como ``uma rigidez cega ao contexto e um viés potencial
baseado no discurso online publicado sobre a situação de rua nas regiões do
mundo com dados super-representados nos LLMs'', em contraste com especialistas
que ajustaram as decisões a realidades codificadas localmente. Trata-se de um
domínio e um tipo de tarefa, que generalizamos para cinco tipos de tarefa
cobrindo recordação, síntese e recomendação. O precedente mais próximo num
domínio regulado é \citet{dahl2024legal}, que fizeram a quatro modelos públicos
perguntas verificáveis sobre casos federais dos Estados Unidos e encontraram
taxas de alucinação entre 58\% e 88\%, mais altas para as cortes inferiores e
as jurisdições menos proeminentes, e mais altas de novo para os casos mais
antigos e os mais recentes, de modo que o conhecimento dos modelos fica atrás
da doutrina em vigor; nosso piso de recordação e a análise da escada na
Seção~\ref{sec:res:t1} são a contraparte ambiental desses dois gradientes. Na
saúde, \citet{kozlakidis2026medical} argumentam que alucinações carregam risco agudo
em ambientes regulatórios com capacidade técnica limitada, onde o descompasso
de idioma e cultura aumenta o erro, um risco que encontra a carga de
mortalidade documentada pelas estimativas do Global Burden of
Disease~\citep{gbd2019risk}.

A questão da implantação já não é hipotética. \citet{kim2026genai} estima
efeitos de produtividade burocrática da IA generativa de forma
quase-experimental, e \citet{robinson2026opensource} documenta que escolher
qual modelo um órgão roda é hoje uma decisão explícita de compra entre opções
de pesos abertos e hospedadas pelo fornecedor. Nossos achados no nível do
modelo falam a essa decisão: os modelos de pesos abertos, como grupo, ficam
$\nhfive$~pontos abaixo dos fechados no composto, ainda que um modelo de pesos
abertos (DeepSeek-V3) lidere a lista, e o modelo português com ajuste regional
seja o mais fraco dos catorze. Dois outros fios enquadram por que uma resposta
errada importa administrativamente. \citet{choroszewicz2026crumple} mostra que
ferramentas de apoio com IA generativa posicionam os profissionais da linha de
frente como ``zonas de deformação moral'', absorvendo a responsabilidade por
falhas originadas a montante, e \citet{cordella2024regulating}, ao analisar a
primeira medida vinculante que um regulador impôs a um serviço de IA generativa
(a intervenção italiana sobre o ChatGPT), argumentam que arcabouços
tecnologicamente neutros ignoram o que é específico desses sistemas e que os
reguladores precisam engajar diretamente seu funcionamento técnico. Sobre a
recepção, \citet{aoki2024explainable} testam se o tipo de explicação muda a
acurácia, a justiça e a confiabilidade percebidas entre os afetados.

Nada disso estabelece se a resposta subjacente está correta, nem se a correção
se distribui de forma uniforme entre os Estados que dependem dela.
Confiabilidade percebida e acurácia medida são grandezas diferentes, e um
sistema pode ser bem explicado, bem governado, bem recebido e errado. Essa é a
lacuna que este estudo enfrenta, num domínio em que a resposta correta está
publicada num registro oficial e pode ser conferida.

\subsection{Posicionamento do presente trabalho}

O desenho difere do trabalho prévio em quatro aspectos: um único domínio
aplicado de alto risco com gabarito oficial verificável, em vez de trivialidades
numéricas ou conhecimento cultural; persona como fator intra-sujeito
pré-especificado, em vez de escolha de enquadramento não controlada; amostragem
de países orientada por teoria em eixos independentes (a divisão
desenvolvido/em desenvolvimento da UNCTAD, recursos linguísticos de Joshi,
desenvolvimento), em vez de conveniência regional; e um teste explícito do
mecanismo de representação no corpus com análise de sensibilidade. Nenhum
estudo anterior combina um domínio regulado de saúde pública, uma manipulação
de persona e um teste de mecanismo pré-especificado.
